\documentclass[11pt]{article}

\usepackage[margin=1in]{geometry}
\usepackage{amsmath}
\usepackage{amssymb}
\usepackage{booktabs}
\usepackage{enumitem}

\title{Tile Packing on a Checkerboard Chip:\\A Mixed-Integer Program}
\author{qSNOW square-packing experiment}
\date{}

\newcommand{\Pset}{\mathcal{P}}
\newcommand{\Vset}{\mathcal{V}}
\newcommand{\Eset}{\mathcal{E}}
\newcommand{\Fp}{F}

\begin{document}
\maketitle

\section{Terminology}

We fix the following vocabulary and use it consistently throughout.
\begin{itemize}[itemsep=2pt]
  \item \textbf{Site} --- an integer coordinate position $(r, c)$ on the chip, where $r$ is
        the row and $c$ is the column. Coordinates are integers with \emph{unit spacing}:
        adjacent lattice lines differ by $1$. Sites live on a checkerboard sublattice: a
        coordinate pair is a site only if it satisfies the parity condition
        $r \equiv c \pmod 2$ (STIM convention; the complementary positions carry ancilla
        measurement hardware and are never tile origins). Consequently, two neighboring sites
        within the same row (or same column) are $2$ coordinate units apart.
  \item \textbf{Physical qubit} --- the data/ancilla qubit hosted at a site. There is exactly
        one physical qubit per site, so ``the physical qubit at site $(r,c)$'' and ``the site
        $(r,c)$'' refer to the same location; we say ``physical qubit'' when emphasizing that
        two tiles must not share hardware, and ``site'' when talking about coordinates.
  \item \textbf{Tile} --- a distance-$d$ rotated-surface-code patch. Placing a tile occupies
        a square block of physical qubits, defined precisely in Section~\ref{sec:geometry}.
  \item \textbf{Origin} --- the site at which a tile is placed, namely the minimum-coordinate
        (top-left) corner of the block of physical qubits it occupies.
\end{itemize}

\section{Problem}

A distance-$d$ tile may be placed with its origin at various sites on the chip. For each
candidate origin we have a Monte-Carlo estimate of the tile's \emph{logical error rate}
(LER) --- the probability that a memory experiment for a tile placed there fails --- obtained
by simulating the tile at that origin and stored in the \texttt{.flake} results file. A
placement is \emph{valid} when its LER does not exceed a fixed threshold. Because each placed
tile occupies a block of physical qubits, two placed tiles may not overlap (share a physical
qubit). We seek to place as many valid, mutually non-overlapping tiles as possible.

\section{Geometry: tiles, footprints, and conflicts}
\label{sec:geometry}

\paragraph{Parameters.}
\begin{itemize}[itemsep=2pt]
  \item $d$ --- the code distance of the tile (e.g.\ $d = 3$).
  \item $s = 2d + 1$ --- the \emph{footprint span}. Here $s$ is a \emph{geometric distance
        measured in coordinate units}, namely the offset from the origin to the far edge of
        the tile along each axis: a tile placed at row $r_p$ occupies rows $r_p$ through
        $r_p + s$. It is \emph{not} a count of occupied sites. Because coordinates have unit
        spacing, this window covers $s+1$ coordinate lines (of which the checkerboard sites
        occupy every other one); $s$ itself is the distance between the first and last line.
        The value $s = 2d+1$ comes from the rotated-surface-code geometry (a distance-$d$
        tile's physical qubits span $2d+1$ coordinate units per side). For $d=3$, $s = 7$
        (an $8$-coordinate-wide window); for $d=5$, $s = 11$.
  \item $\tau$ --- the LER validity threshold (default $\tau = 0.20$).
\end{itemize}

\paragraph{Footprint.}
A tile placed at origin $p = (r_p, c_p)$ --- where $r_p$ and $c_p$ are the row and column of
that origin --- occupies every site inside the axis-aligned coordinate window that extends
$s$ units from the origin toward increasing row and column. We call this set of occupied
sites the tile's \emph{footprint}:
\[
  \Fp(p) \;=\; \bigl\{\, (r, c) \;:\; r_p \le r \le r_p + s,\;\;
  c_p \le c \le c_p + s,\;\; r \equiv c \!\!\pmod 2 \,\bigr\}.
\]
Thus the footprint is a square block of physical qubits: its coordinate window runs from
$r_p$ to $r_p + s$ in the row direction and from $c_p$ to $c_p + s$ in the column direction
(inclusive). The far edge is therefore a distance $s$ (in coordinate units) from the origin,
and the window covers $s+1$ coordinate lines on each side; the checkerboard sites occupy
every other position inside it. For $d=3$ ($s=7$) this is an $8$-coordinate-wide window
anchored at the origin.

\paragraph{Conflict.}
Two distinct placements \emph{conflict} when their footprints share at least one site (i.e.\
the two tiles would compete for the same physical qubit). Since a footprint is an
axis-aligned window of fixed side $s$, the row windows $[r_p, r_p + s]$ and $[r_q, r_q + s]$
of two tiles overlap exactly when $|r_p - r_q| \le s$, and likewise for columns. Two footprints
(2-D windows) share a site precisely when they overlap in \emph{both} axes, giving the
conflict test
\[
  p \sim q
  \quad\Longleftrightarrow\quad
  p \neq q,\;\;
  |r_p - r_q| \le s
  \;\;\text{and}\;\;
  |c_p - c_q| \le s .
\]

\paragraph{Worked example ($d=3$, so $s=7$).}
A tile at origin $(0,0)$ occupies the sites in $[0,7]\times[0,7]$. A second tile at $(0,8)$
does \emph{not} conflict: the columns differ by $8 > 7$, so the two column windows are
disjoint. A tile at $(1,7)$ \emph{does} conflict: both coordinate differences ($1$ and $7$)
are $\le 7$, so the footprints share physical qubits. This matches the blocking pattern
observed directly in the simulator.

\section{Data and sets}

\begin{itemize}[itemsep=2pt]
  \item $\Pset$ --- the set of candidate origins read from the results file. Each origin
        $p \in \Pset$ is a site $(r_p, c_p)$ (so $r_p \equiv c_p \pmod 2$) and carries a
        measured LER $\ell_p \in [0,1]$.
  \item Valid placements:
        \[
          \Vset \;=\; \{\, p \in \Pset \;:\; \ell_p \le \tau \,\},
        \]
        the origins whose tile is good enough to use at threshold $\tau$.
  \item Conflict edges among valid placements:
        \[
          \Eset \;=\; \bigl\{\, \{p,q\} \;:\; p,q \in \Vset,\; p \sim q \,\bigr\},
        \]
        using the conflict relation $\sim$ from Section~\ref{sec:geometry}.
\end{itemize}

\section{Decision variables}
For each valid origin we introduce one binary variable:
\[
  x_p \in \{0,1\}, \qquad p \in \Vset,
  \qquad
  x_p =
  \begin{cases}
    1 & \text{if a tile is placed at origin } p,\\
    0 & \text{otherwise.}
  \end{cases}
\]

\section{Model}

\begin{align}
  \max_{x} \quad & \sum_{p \in \Vset} x_p
    && \text{(maximize the number of placed tiles)} \label{eq:obj}\\[4pt]
  \text{s.t.} \quad
    & x_p + x_q \le 1
    && \forall\, \{p,q\} \in \Eset
    && \text{(no two placed tiles overlap)} \label{eq:pair}\\[2pt]
    & x_p \in \{0,1\}
    && \forall\, p \in \Vset. && \label{eq:bin}
\end{align}

Objective~\eqref{eq:obj} counts placed tiles; constraint~\eqref{eq:pair} forbids placing
both members of any conflicting pair. This is a maximum-cardinality packing, equivalently a
maximum independent set on the conflict graph $(\Vset, \Eset)$.

\section{Tighter (clique) formulation}

The pairwise constraints \eqref{eq:pair} are correct but yield a weak linear relaxation
(the LP can set every $x_p = \tfrac12$). We can tighten the model as follows.

Because footprints are axis-aligned windows, any collection of tiles that pairwise conflict
has a common site: for axis-aligned boxes, pairwise intersection implies a common point
(Helly's theorem in the plane). Hence every maximal clique of the conflict graph corresponds
to a single site that lies in the footprint of all the tiles in the clique. This lets us
replace the many edge constraints by one constraint per site.

For a site $u = (u_r, u_c)$ --- where $u_r$ and $u_c$ are its row and column --- let
\[
  \mathcal{C}_u \;=\; \bigl\{\, p \in \Vset \;:\; u \in \Fp(p) \,\bigr\}
  \;=\; \bigl\{\, p \in \Vset \;:\; r_p \le u_r \le r_p + s
  \;\text{and}\; c_p \le u_c \le c_p + s \,\bigr\}
\]
be the set of valid placements whose footprint covers site $u$. Any two placements in
$\mathcal{C}_u$ share $u$ and therefore conflict, so at most one of them may be used:
\begin{equation}
  \sum_{p \in \mathcal{C}_u} x_p \le 1
  \qquad \forall\, \text{sites } u .
  \label{eq:clique}
\end{equation}
Each such constraint dominates all the pairwise edges among the tiles covering $u$, so
\eqref{eq:clique} gives a stronger LP relaxation than \eqref{eq:pair} while describing the
same set of integer solutions. In the solver we enumerate only the \emph{maximal} covered
windows (equivalently, distinct sets $\mathcal{C}_u$ that are not contained in another),
which suffices to cover every conflicting pair.

\section{Notes}
\begin{itemize}[itemsep=2pt]
  \item The objective can be generalized to a weighted packing
        $\max \sum_p w_p x_p$ (e.g.\ $w_p = 1 - \ell_p$ to prefer lower-error placements);
        the constraints are unchanged.
  \item Chip-boundary feasibility is already baked into $\Pset$: the results file only
        contains origins whose full footprint fits on the chip.
\end{itemize}

\end{document}
